\documentclass[10pt, oneside]{article}

\usepackage{amsmath}
\usepackage{amsthm}
\usepackage{amssymb}
\usepackage{calrsfs}
\usepackage{wasysym}
\usepackage{verbatim}
\usepackage{bbm}
\usepackage{color}
\usepackage{graphics}
\usepackage{geometry}
\usepackage{graphicx}

\geometry{tmargin=.75in, bmargin=.75in, lmargin=.75in, rmargin = .75in}  

\newcommand{\R}{\mathbb{R}}
\newcommand{\C}{\mathbb{C}}
\newcommand{\Z}{\mathbb{Z}}
\newcommand{\N}{\mathbb{N}}
\newcommand{\Q}{\mathbb{Q}}
\newcommand{\Cdot}{\boldsymbol{\cdot}}

\newtheorem{thm}{Theorem}
\newtheorem{defn}{Definition}
\newtheorem{conv}{Convention}
\newtheorem{rem}{Remark}
\newtheorem{lem}{Lemma}
\newtheorem{cor}{Corollary}

\title{System-2 Course: Session [Session Number]}

\author{[Full Name] - [Student Number]}
\date{Spring 2025}

\begin{document}

\maketitle

\tableofcontents

\pagebreak

\section{Section 1}

\subsection{Subsection 1}

Text ...

\begin{itemize}
    \item Item 1
    \item Item 2
\end{itemize}

\begin{rem}
    Text ...
\end{rem}

\begin{defn}
    Text ...
\end{defn}

\subsection{Subsection 2}

\begin{equation}
    E = m \times c^2
\end{equation}

\begin{thm}
    Text ...  
\end{thm}

\begin{lem}
    Text ...  
\end{lem}

\section{Section 2}

% \begin{figure}
%     \centering
%     \includegraphics[width=0.5\linewidth]{image_address}
%     \caption{Caption}
%     \label{fig:enter-label}
% \end{figure}

\end{document}